\begin{theorem}[Resolving the partner label in a regional charge pair]
    \label{thm:peps_closed_charge_pair_resolution}
    \lean{TNLean.PEPS.sum_openRegionWeight_regularCorrelatedChargePair}
    \leanok
    Let $R$ be a region of a finite simple graph, with finite group $G$ on its
    bonds. Let $e_0,e_1$ be two internal bonds, $u$ an arbitrary bond assignment,
    and $\theta$ a fixed crossing configuration. Resolving the label on $e_1$
    by the diagonal indicators $\mathbf{1}_{\eta_{e_1}=q}$ gives
    \begin{align}
        \sum_{q\in G} B_R^{u,\,\chi(pq^{-1}\eta_{e_0})
                \mathbf{1}_{\eta_{e_1}=q}}(\theta)
         & =B_R^{u,\,\chi(p\eta_{e_1}^{-1}\eta_{e_0})}(\theta).\notag
    \end{align}
    Both sides are the actual open contractions of the original site tensors.
    The character weight remains correlated between the two labels.
    This is the finite-contraction identity underlying the creation of charge
    pairs in \cite[Section~6.6]{Schuch2010PEPS}, source lines 2505--2535.
    \uses{def:peps_regular_correlated_charge_state}
\end{theorem}
\begin{proof}\leanok
    Expand the original open contraction and interchange its finite label sums.
    The diagonal indicator fixes $q=\eta_{e_1}$.
\end{proof}

\begin{theorem}[Original-spin preparation of the closed correlated pair]
    \label{thm:peps_closed_charge_pair_preparation}
    \lean{TNLean.PEPS.exists_unitary_regularClosedChargePairPreparation,
        TNLean.PEPS.regularClosedChargePairPreparation_normSq_ne_zero}
    \leanok
    Let the site tensors of a finite simple graph be regular $G$-isometric.
    Let $R$ have a specified spanning tree and root, and let $e_0,e_1$ be two
    distinct internal bonds. For every irreducible character $\chi$ occurring in
    the regular representation and every $p\in G$, there is a unitary $W$ on the
    original spins in $R$ such that
    \begin{align}
        (W\otimes I_{R^c})\Psi_0 & =\Psi_{\chi,p},\notag \\
        \Psi_{\chi,p}(\sigma)    & =\sum_{\eta:E\to G}
        \chi(p\eta_{e_1}^{-1}\eta_{e_0})
        \prod_v a_v(\eta|_v,\sigma_v).\notag
    \end{align}
    Here $\Psi_0$ is the actual closed contraction without insertions.
    Consequently $\Psi_{\chi,p}\ne0$ and
    $\langle\Psi_{\chi,p},\Psi_{\chi,p}\rangle
        =\langle\Psi_0,\Psi_0\rangle$.
    The local operation is uniform in every crossing configuration.

    This auxiliary finite-region consequence of SCP10, lines 2505--2558,
    retains the supplied region, spanning tree and two distinct internal bonds.
    It does not identify the prescribed charge--flux braid, its flux-string
    background, or the reunion experiment of lines 2560--2615; see
    \cite{gap:rmp_peps_quantum_double_g_isometry}.
    \uses{thm:peps_closed_charge_pair_resolution,
        thm:peps_regular_correlated_charge_cut}
\end{theorem}
\begin{proof}\leanok
    Local $G$-isometry supplies the original-spin unitary preparing each correlated
    open column. Resolve the partner label and contract the unchanged exterior
    columns. The resulting closed identity uses the same unitary. Unitarity
    preserves the inner product, and the unmodified regular $G$-injective closed
    contraction is nonzero.
\end{proof}
